% Generated by scripts/materialize_unified_paper.py; do not edit.
\section{Ramsey rates and imported interfaces}\label{up:sec:framework}

\subsection{Symmetric homogeneous rate}

Let \(U:[0,1]\to\R\) be continuous with \(U(0)=0\).  Define its symmetric
homogeneous extension by
\[
 \wideU(x,y)=
 \begin{cases}
 \max\{x,y\}\,U\!\left(\dfrac{\min\{x,y\}}{\max\{x,y\}}\right),
     &\max\{x,y\}>0,\\[2mm]
 0,&x=y=0.
 \end{cases}
 \tag{U2.1}\label{up:main:eq:2-1}
\]
For residual problems only, we extend the notation by
\(R(a,b)=1\) when \(\min\{a,b\}\le0\); this means that no further page
vertices are needed in at least one colour.  It is not a new convention for
positive Ramsey numbers.

\begin{lemma}[Uniform homogeneous extension]\label{up:lem:homogeneous}
Assume in addition that \(U\ge0\).  Suppose \eqref{up:main:eq:1-1} holds uniformly for
\(1\le\ell\le k\), and the same bound is
available after exchanging the two colours.  Then, uniformly for
\(x,y\in[0,1]\),
\[
 \log R(\lfloor xk\rfloor,\lfloor yk\rfloor)
 \le k\wideU(x,y)+o(k).
 \tag{U2.2}\label{up:main:eq:2-2}
\]
\end{lemma}

\begin{proof}
The extension \(\wideU\) is continuous on the square (continuity at the
origin uses \(U(0)=0\)), hence uniformly continuous.  Fix
\(\epsilon>0\), and choose \(K_0\) from the uniform epsilon formulation of
\eqref{up:main:eq:1-1}.  For \(x\ge y\), put \(k'=\lfloor xk\rfloor\) and
\(\ell'=\lfloor yk\rfloor\).  If \(k'\ge K_0\) and \(\ell'\ge1\), then
\[
 \log R(k',\ell')\le
 k' U(\ell'/k')+\epsilon k'
 =k\wideU(k'/k,\ell'/k)+\epsilon k.
\]
Uniform continuity and the \(O(k^{-1})\) perturbation of the normalized
coordinates replace the last \(\wideU\) by \(\wideU(x,y)+o(1)\), uniformly.
If \(k'<K_0\), only finitely many positive Ramsey numbers occur, so their
logarithms are \(O_{K_0}(1)=o(k)\); nonnegativity of \(\wideU\) permits the
same upper bound.  If one floored coordinate is zero, the residual convention
gives logarithmic cost zero.  The case \(y\ge x\) follows by symmetry.  Since
\(\epsilon\) is arbitrary, these cases give one uniform \(o(k)\).
\end{proof}

\subsection{Yang--Mao interfaces}

The following package is extracted from the v1 source of
Yang--Mao \cite{YangMao2026}.  Its source member and line map are frozen in
\cref{up:app:upper-certificate}.  We state the complete specialization used
later, so that the conditional theorem is falsifiable without guessing what
the words ``book interface'' mean.

For completeness, we fix the precise meaning of the correlation input.  We
write \(\cG_2^{(3)}(\beta,C)\) when the following assertion holds: for every
finite set \(\mathcal X\), independent identically distributed
\(U,U'\in\mathcal X\), real Hilbert spaces \(\mathcal H_1,\mathcal H_2\),
and maps \(\sigma_i:\mathcal X\to\mathcal H_i\), if
\(Z_i=\langle\sigma_i(U),\sigma_i(U')\rangle\), then there are
\(i\in\{1,2\}\) and \(\lambda\ge-1\) such that
\[
 \mathbb P\bigl(Z_i\ge\lambda,\ Z_{3-i}\ge-1\bigr)
 \ge \beta\exp\!\left(-C(\lambda+1)^{1/3}\right).
 \tag{U2.2a}\label{up:eq:correlation-definition}
\]

For \(0<\beta\le1\), \(C>0\), and parameters
\(p,\delta,\lambda_0>0\), put
\[
 L_\delta=\log(1/\delta),\qquad L_p=\log(1/p),\qquad
 \rho_\beta=\log(2/\beta),
\]
\[
 \Pi=\frac{3\delta}{p}+\frac{6L_\delta L_p}{\lambda_0},\qquad
 q=L_p+\Pi,
 \tag{U2.3}\label{up:eq:book-costs}
\]
and
\[
 \Xi=2\rho_\beta+4C\lambda_0^{1/3}
       +\frac{12CL_\delta}{\lambda_0^{2/3}}.
 \tag{U2.4}\label{up:eq:reservoir-cost}
\]

\begin{assumption}[Retained-spine source interface]\label{up:ass:yang-mao}
Fix \(0<\eta\), \(0<p<1/2-\eta\),
\(0<\delta\le\min\{p/4,1/4\}\),
\(\lambda_0\ge\max\{2,6\log(1/\delta)\}\), and \(\tau>0\).
Suppose the correlation property \(\cG_2^{(3)}(\beta,C)\) holds.  The
version-pinned Yang--Mao specialization supplies the following conclusions
for all sufficiently large \(k\).
\begin{enumerate}[label=(\alph*)]
\item For every positive integer \(N\), every red--blue colouring of
      \(K_N\) admits pairwise disjoint sets \(S_R,S_B,W\) such that
      \(S_R\) is a red clique, \(S_B\) is a blue clique, every edge from
      \(S_R\) to \(W\) is red, and every edge from \(S_B\) to \(W\) is
      blue.  Writing \(s_i=|S_i|\), one has
      \[
      |W|\ge \left(\frac{1+\eta}{2}\right)^{s_R+s_B}N
      \tag{U2.5}\label{up:main:eq:2-5}
      \]
      and every \(w\in W\) has at least
      \((1/2-\eta)|W|-1\) red neighbours and at least that many blue
      neighbours inside \(W\).
\item Let \(t,m\) be positive integers with
      \(t\ge\lambda_0\delta^{-1/2}\).  If every \(w\in W\) has at least
      \(p|W|\) neighbours of each colour inside \(W\), and
      \[
       |W|\ge p^{-t}e^{\Pi t}m,
       \qquad |W|\ge4t\,e^{2t\Xi},
       \tag{U2.6}\label{up:main:eq:2-6}
      \]
      then the parameterized-book theorem, applied with
      \(X=Y_R=Y_B=W\), returns a colour \(i\in\{R,B\}\), a monochromatic
      \(t\)-clique \(T\subseteq W\) of colour \(i\), and
      \(P_0\subseteq W\setminus T\), \(|P_0|\ge m\), such that every edge
      from \(T\) to \(P_0\) has colour \(i\).
\item The compatibility conclusion is the following exact residual
      statement.  If \(i=R\), then a red
      \(K_{k-s_R-t}\) in \(P_0\) (when this order is positive) extends with
      \(S_R\cup T\), while a blue \(K_{k-s_B}\) extends with \(S_B\).
      If \(i=B\), interchange the colours, so the two residual orders are
      \(k-s_R\) and \(k-s_B-t\).  A nonpositive residual order means that
      the retained spines already contain the required \(K_k\).
\end{enumerate}
For the proof below the exact correlation input is
\[
 \beta=\frac{330867}{500000},\qquad
 C=\frac{12366348252219}{1250000000000}.
 \tag{U2.7}\label{up:eq:correlation-input}
\]
\end{assumption}

The point of stating \cref{up:ass:yang-mao} is auditability.  The correlation
property in \eqref{up:eq:correlation-input} is proved in this paper; the regularization and book
compatibility are imported.  The source line mapping and pinned archive hash
are included in the reproducibility record.

\subsection{A general retained-spine transfer}

Put
\[
 c_\eta=\log\frac{2}{1+\eta}.
\]
For a normalized spine vector \(\sigma=(\sigma_R,\sigma_B)\in[0,1]^2\),
define
\[
 D(\sigma)=c_\eta(\sigma_R+\sigma_B)
 +\wideU(1-\sigma_R,1-\sigma_B).
 \tag{U2.8}\label{up:main:eq:2-8}
\]
Here \(q\) and \(\Xi\) are the explicit quantities in \eqref{up:eq:book-costs}--\eqref{up:eq:reservoir-cost}.  Define
\begin{align}
 M_\tau(\sigma)=\max\{&
 \wideU((1-\sigma_R-\tau)_+,1-\sigma_B),\nonumber\\
 &\wideU(1-\sigma_R,(1-\sigma_B-\tau)_+)\},\tag{U2.9}\label{up:main:eq:2-9}\\
 P(\sigma)&=c_\eta(\sigma_R+\sigma_B)+\tau q+M_\tau(\sigma),\tag{U2.10}\label{up:main:eq:2-10}\\
 Q(\sigma)&=c_\eta(\sigma_R+\sigma_B)+2\tau\Xi,\tag{U2.11}\label{up:main:eq:2-11}\\
 B(\sigma)&=\max\{P(\sigma),Q(\sigma)\}.\tag{U2.12}\label{up:main:eq:2-12}
\end{align}
Finally set
\[
 C_*=\max_{\sigma\in[0,1]^2}\min\{D(\sigma),B(\sigma)\}.
 \tag{U2.13}\label{up:main:eq:2-13}
\]

\begin{theorem}[Retained-spine transfer]\label{up:thm:transfer-general}
Under \cref{up:ass:yang-mao} and the uniform rate \eqref{up:main:eq:2-2},
\[
 R(k,k)\le \exp(C_*k+o(k)).
\]
\end{theorem}

\begin{proof}
Let \(N=\lceil e^{(C_*+\varepsilon)k}\rceil\) and apply regularization.
If either preliminary spine already has order at least \(k\), there is
nothing to prove.  Otherwise write \(s_i=\sigma_i k+O(1)\).  From \eqref{up:main:eq:2-5},
\[
 \log|W|\ge(C_*+\varepsilon)k
 -c_\eta(s_R+s_B)+o(k).
 \tag{U2.14}\label{up:main:eq:2-14}
\]
At this \(\sigma\), choose the smaller of the direct and book branches.

If \(D(\sigma)\le B(\sigma)\), then \(D(\sigma)\le C_*\).  Equations
\eqref{up:main:eq:2-2}, \eqref{up:main:eq:2-8} and \eqref{up:main:eq:2-14} imply
\(|W|\ge R(k-s_R,k-s_B)\).  The resulting red or blue clique extends with
the corresponding preliminary spine.

If \(B(\sigma)<D(\sigma)\), then both \(P(\sigma)\le C_*\) and
\(Q(\sigma)\le C_*\).  The former supplies the page size needed for either
possible book colour, using \eqref{up:main:eq:2-2} and \eqref{up:main:eq:2-9}; the latter supplies the common
reservoir threshold.  More explicitly, with \(t=\lfloor\tau k\rfloor\)
and \(m\) the larger of the two residual Ramsey thresholds, the two strict
exponent inequalities give exactly \eqref{up:main:eq:2-6}.  Since \(p<1/2-\eta\) and the
reservoir is exponential, the additive \(-1\) in the degree conclusion is
absorbed; fixed \(\tau>0\) also gives
\(t\ge\lambda_0\delta^{-1/2}\).  The parameterized-book interface produces
the \(t\)-spine and a page set of size at least \(m\), and compatibility
extends the appropriate page clique with the retained spines to a
monochromatic \(K_k\).  The fixed \(\varepsilon\)-slack absorbs every
uniform \(o(k)\) term.
\end{proof}
